\section{What the conditional chain composes}
The chain uses a common topological and measured substrate, its sheaf of
global sections, a supplied complete metric on those sections, a neural field,
a thermodynamic cover $T$, and a reflexive boundary with prediction map $P$.
The precise type and instance parameters are printed in the appendix.

\subsection{The common downstream argument}
Suppose a sequence of energies coarse-grains to $L$, the elapsed parameter
$\tau$ is positive, and the following named implications hold:
\begin{itemize}
\item $E56$: coarse-graining implies $D>0$, $K\geq0$, $L=K$, and an
identification of the supplied field with the parameters $K,D$.
\item $E67$: those scalar conditions imply $K>2D$.
\item $E78$: a coherent order parameter implies that the supplied cover is
reached by the specified relaxation at a coupling at least $K$.
\item $E89$: that reachability supplies a section $s$ which restricts to the
local data of $T$, is fixed by $P$, and comes with the stated Lipschitz bound
on $P$ at \texttt{resonanceRate}$(K,D,\tau)$.
\end{itemize}

\begin{claim}
There is a global section compatible with a thermodynamic cover which is the
unique fixed point of the prediction map.
\end{claim}
\begin{proof}
Apply $E56$ to the coarse-graining hypothesis and $E67$ to its scalar
conclusion. The supercritical existence theorem supplies a coherent order
parameter. Apply $E78$ and then $E89$ to obtain the section $s$, its local
restrictions, the Lipschitz bound, and $P(s)=s$.

The section $s$ makes the space of global sections nonempty. The theorem
\nolinkurl{self_of_coherent_order_parameter}, using completeness, the
Lipschitz bound, coherence, and $\tau>0$, applies the contraction argument to
obtain a unique fixed point $p$. Since $s$ is fixed, uniqueness gives $s=p$.
Every fixed point $t$ equals $p$ and hence $s$. Pair $s$ with the supplied
cover $T$ and its known restriction equations. This is exactly the predicate
\texttt{UnifiedSelf}.
\end{proof}
\formal{downstream-chain}
\scope The identity between the glued section and a fixed section is supplied
in $E89$; the contraction argument proves uniqueness. The concrete-field
conjunct of $E56$ is retained as an input constraint but is not used by this
proof's downstream fixed-point argument. The source also establishes a local
\texttt{Unity} fact which the concluding construction does not reference.
Its constants can still occur in the unnormalised proof term; a syntactic
dependency list alone therefore does not identify minimal premises.

\subsection{The passive composition}
The passive chain additionally supplies a finite nonempty register, its
statistical-mechanics instance and update, a field and vacuum set, and measurable
carrier types for predictive information. Its remaining named implications are:
\begin{itemize}
\item $E12$: the capacity bound implies that the field leaves the vacuum set.
\item $E23$: leaving that set implies the register update is not surjective.
\item $E34$: dissipation supplies a predictive-dissipation structure whose
thermal scale and dissipated work match the register's temperature and heat,
and whose nonpredictive information is strictly positive.
\item $E45$: the predictive bound implies the chosen coarse-graining statement.
\end{itemize}

\begin{claim}
Under $E12$ through $E89$ and the stated structural hypotheses, the passive
chain concludes \texttt{UnifiedSelf} for the supplied reflexive boundary.
\end{claim}
\begin{proof}
Finiteness gives the capacity bound. Apply $E12$ to obtain vacuum departure,
then $E23$ to obtain a nonsurjective update. The proved dissipation theorem,
using the supplied statistical-mechanics structure, gives positive heat.
Apply $E34$ and retain its predictive-dissipation witness $R$. The bound
associated with $R$ gives the predictive node. Apply $E45$ to obtain
coarse-graining and then invoke the downstream claim with $E56$ through $E89$.
\end{proof}
\formal{passive-chain}
\scope The eight named implications remain hypotheses. Conditions inside
the statistical-mechanics, thermodynamic-cover, and other structures are
additional obligations. The equalities and positivity in $E34$ constrain
which witness is admitted, while the passage through $E45$ uses its
predictive structure. The conclusion has no predicate asserting experience
or biological implementation.
